% !TeX program = xelatex
% CUMCM 2026 C 题《AI 工具使用详情》——据实填写，提交前导出 PDF
\documentclass[a4paper,11pt]{ctexart}
\usepackage[margin=25mm]{geometry}
\usepackage{booktabs,tabularx,array,longtable,enumitem,hyperref}
\hypersetup{hidelinks}
\setlength{\parindent}{2em}
\renewcommand{\arraystretch}{1.35}
\title{AI 工具使用详情}
\author{CUMCM 2026 参赛队（C 题：微网与外部电网电力调控策略）}
\date{2026 年 9 月}
\begin{document}
\maketitle

\section*{一、所用 AI 工具名称、版本或型号}
\begin{tabularx}{\textwidth}{>{\bfseries}p{0.20\textwidth} p{0.24\textwidth} p{0.16\textwidth} X}
\toprule
工具名称 & 版本／型号 & 开发机构 & 使用日期与环节 \\
\midrule
OpenCode（AI 编程智能体） & OpenCode 1.18.30；调用模型：DeepSeek V4 Flash & OpenCode 开源生态；模型由 DeepSeek 提供 & 2026-09-09 至 2026-09-13：建模辅助、程序编写与调试、数值实验、图表绘制 \\
ChatGPT & GPT-5.6 Thinking & OpenAI & 2026-08-03 起：论文 LaTeX 排版与格式规范；2026-09：求解结果的独立平行审计与稳健性对照 \\
\bottomrule
\end{tabularx}

\section*{二、具体使用目的和环节}
本队将 AI 定位为\textbf{放大器而非代理大脑}：AI 用于放大参赛队的思考与执行效率，
\textbf{核心建模与分析由参赛队主导}，并对 AI 参与完成的内容逐项人工审查与核实；
\textbf{代码虽由 AI 辅助编写，但由参赛队对代码负责并承担 review 责任}。具体环节如下：

\begin{enumerate}[leftmargin=2.5em,itemsep=2pt]
\item \textbf{建模思考与方案对照}：提供备选建模路径及优缺点对照（如问题二预测目标的选择、
问题三调整机制的取舍、问题四电价信息口径 H/G 的对比），供参赛队评估后决策；
\item \textbf{模型设计与公式整理}：将参赛队确定的建模思路整理为规范公式、约束与算法流程
（组合预测、按下游费用标定的权重、分位数裕度修正、三层结算 LP、跨日滚动结构等），并逐项复核；
\item \textbf{程序编写、数值实验与调试}：在参赛队主导设计下，AI 辅助完成求解与实验程序的
编写和调试，批量执行参数搜索、灵敏度分析、蒙特卡洛与对照实验，生成对照表供参赛队比较与取舍；
\item \textbf{初步数据分析与绘图}：快速完成数据概览、周期结构与误差分布等初步分析，并快速
出图；
\item \textbf{论文写作与排版}：按 CUMCM 模板整理 LaTeX 工程，排查并修复 XeLaTeX 编译问题，
按官方格式规范调整版式；
\item \textbf{验证与审计}：执行五份结果文件结构审计、写作门禁检查与无前视审计；并由 ChatGPT
对 Q2/Q3/Q4 的求解结果做独立平行复算与稳健性对照，用于交叉验证。
\end{enumerate}

\section*{三、主要提示方式与使用过程（典型示例）}
\textbf{主要提示方式}：先由参赛队给出建模目标、假设、约束与验收标准，再要求 AI 提供可运行的
实现与可复核的对照数据；对 AI 输出\textbf{先核验、后采纳}，关键结论一律回到程序输出核对。
典型示例如下。

\begin{longtable}{>{\centering\arraybackslash}p{0.05\textwidth}p{0.38\textwidth}p{0.47\textwidth}}
\toprule
序号 & 请求／过程 & 输出与采纳、人工修改情况 \\
\midrule
\endfirsthead
\toprule
序号 & 请求／过程 & 输出与采纳、人工修改情况 \\
\midrule
\endhead
1 & 对已完成的求解程序做执行层独立审计（ChatGPT 平行验证，2026-09） & 指出“全天事后 LP 存在前视，会系统性低估紧急购电”，参赛队核查后改为逐槽因果执行，全部正式数字重算。\textbf{采纳}。 \\
2 & 三源凸组合预测与“按下游费用标定权重”由参赛队提出思路，AI 负责实现与数值执行（OpenCode，2026-09） & 实现三源凸组合与按历史费用逐日标定的 EWMA 权重，并生成全年对照数据。\textbf{采纳}（人的灵感、AI 执行）。 \\
3 & 具有机器学习背景的队员提出在网格搜索上做梯度精细优化并借鉴 Adam 自适应步长，交由 AI 实现（OpenCode，2026-09） & 在 softmax 参数上实现“网格定位 + Adam 中心差分梯度精化”，给出对照数据：精化带来显著改进。\textbf{采纳}。 \\
4 & 对不熟悉的模型评价方法，先向 AI 请教原理与适用条件（OpenCode/ChatGPT，2026-09） & 讲解指标含义与适用前提；参赛队学习并确认理解后，才用于本问的检验设计。\textbf{先学后用}。 \\
5 & 参数平滑系数先按 Adam 一阶动量经验取 0.9，请评估更合适的取值（OpenCode，2026-09） & 建议学习指数加权平均并用数据检验平滑强度，最终改用 0.1 的新旧混合系数（电价权重动态修正）。\textbf{采纳}。 \\
6 & 结构审计脚本与论文排版规范整理（OpenCode/ChatGPT，2026-09） & 生成结果文件结构审计脚本（一次运行全部通过）；按官方规范整理 LaTeX 版式并修复编译问题。\textbf{采纳}。 \\
\bottomrule
\end{longtable}

\section*{四、对 AI 输出的采纳、人工修改和核验的主要情况}
\begin{itemize}[leftmargin=2em,itemsep=2pt]
\item \textbf{完全采纳}：因果执行修正（消除前视，影响全部正式数字）、问题二调优配置、
五份结果文件结构审计脚本、网格+梯度精化与参数平滑定值、论文排版与格式规范；
\item \textbf{部分采纳（经人工修改）}：AI 输出的公式推导、文字草稿与图表均经参赛队复核修改，
包括摘要精简、数值换算与口径统一、图表命名与说明、参考文献逐条核实；
\item \textbf{未采纳／暂缓}：回归参考量方案（净收益约 0.24\%）、预测层缩放/校准类修饰
（经费模型吸收、费用不变）、深度学习第三预测源（单年数据下季节风险大）——均留有实验记录；
\item \textbf{代码责任与可复核性}：代码虽由 AI 辅助编写，但由参赛队对代码负责并承担 review
责任；为便于复核，参赛队主动要求 AI 编写函数级与代码块级注释，并在关键位置标注数学公式
（LP 目标与约束、结算口径、采样与相位映射）；关键算法代码由参赛队逐行阅读复核，并配合回归
测试与五份结果文件审计，通过后才进入正式流程；
\item \textbf{参赛队负责}：赛题理解、建模主线与信息口径选择、参数与方案拍板、结果核验、
论文定稿与提交。全部数值结果由固定随机种子的程序重新生成并逐项校验。
\end{itemize}

\section*{五、完整性声明}
本文件覆盖本队可核查的 AI 使用记录（截至 2026-09-13），与论文的《AI 工具使用声明》及
支撑材料一一对应；工具与版本以使用当时为准。如提交前另有补充使用，将更新本文件并重新导出 PDF。

\end{document}
